% TC-0008 test-case definition (WSP test-case template).
\documentclass[letterpaper]{article}

\usepackage[letterpaper,margin=1in]{geometry}

% Compile from the repository root so this path resolves to the pinned WSP
% submodule.
\input{wsp/testing/test-case-library.tex}

\begin{document}

\def\TCID{TC-0008}
\def\TCTitle{Corpus 212: FlatScrollBar defaults, range, events and scripted scrolling}
\def\TCPurpose{%
  Verify the form-file values and the range of two FlatScrollBar controls, error 380 for a Value outside the range, Change and Scroll, and the scripted scroll and scrollby actions.}
\def\TCPriority{\TCPriorityHigh}

\def\TCRequirementRef{REQ-0002}
\def\TCDesignRef{Not applicable.}
\def\TCFeatureRef{\texttt{tests/corpus/212-mscomctl2-flatscrollbar.vbp}, \texttt{.frm}, \texttt{.ui}, \texttt{.debug} and \texttt{.expected}.}
\def\TCTestTechnique{%
  Scenario test: a program that prints what it observes, compared with a stored transcript.}

\def\TCPreconditions{%
  A \texttt{wfc} executable that includes this library (WFC version 0.8.28 or later) and the files of the case in \texttt{tests/corpus}.}
\def\TCEnvironment{%
  Windows x64. The \texttt{wfc} executable of WFC, given to CMake as \texttt{WFC\_EXE}; the program is run by the interpreter (\texttt{--interpret}) with the system common controls. The case is registered only when \texttt{WFC\_EXE} is set; without it CTest reports the case as skipped.}
\def\TCAssumptions{%
  Not applicable.}

\def\TCInputData{%
  Value 500 and 50; the script lines \texttt{scroll Flat1 70}, \texttt{scrollby Flat1 -5}, \texttt{scroll Flat2 99999} and \texttt{scroll Flat1 70}.}
\def\TCInitialState{%
  No output of an earlier run is used. CTest test name: \texttt{TC-0008-corpus-flatscrollbar}.}

\def\TCProcedure{%
  \begin{enumerate}
    \item Run \texttt{wfc --interpret tests/corpus/212-mscomctl2-flatscrollbar.vbp}. Form\_Load prints the start values, sets Value 500 (error) and Value 50, and prints the type name and the default property.
    \item The script drags Flat1 to 70, scrolls it by -5, drags Flat2 to 99999 and Flat1 to 70 again, then closes the form.
    \item Compare standard output (the form's Print) with \texttt{212-mscomctl2-flatscrollbar.expected}; there is no Immediate-window file for this case.
  \end{enumerate}}
\def\TCExpectedResult{%
  The start line shows the values of the form file (Value 20, Min 10, Max 100 for Flat1, Max 32767 for Flat2, LargeChange 10, SmallChange 1); Value 500 gives \texttt{380 Invalid property value}; setting 50 prints \texttt{change 50}; a drag to 70 prints \texttt{scroll 70} and \texttt{change 70}, a scroll by -5 prints only \texttt{change 65}, and Flat2 stops at 32767 (\texttt{change2 32767}). The output equals the stored file.}
\def\TCPostConditions{%
  No files are written; the interpreter's windows are closed by the script.}

\TCTemplate

\end{document}
